\item Un medidor de presión en un cilindro registra la presión manométrica, que es la diferencia entre las presiones interior y exterior $P_0$. La presión manométrica se denota por $P_g$. Cuando el tanque está lleno, la masa de gas en él es $m_i$ a la presión atmosférica $P_0$. Suponga que la temperatura del cilindro permanece constante, entonces demuestre que la masa del gas \textit{restante} en el cilindro cuando la lectura de la presión es $P_{gf}$ está dada por
$$m_f = m_i \left( \frac{P_{gf} + P_0}{P_{gi} + P_0} \right)$$
